% Section: first-order-definability

\documentclass[../../../include/open-logic-section]{subfiles}

\begin{document}

\olfileid[fr]{nml}{frd}{fol}

\olsection{Définissabilité au premier ordre}

Nous avons vu que certaines propriétés des relations
d'accessibilité des cadres peuvent être définies par des
!!{formula}s modales. La symétrie des cadres, par exemple, est
définie par la !!{formula}~\Ax{B},
$p \lif \Box\Diamond p$. Toutes les conditions rencontrées
jusqu'ici s'expriment aussi par des !!{formula}s du premier
ordre dans un !!{language} ne comportant qu'un seul
!!{predicate} binaire. Ainsi, la symétrie est définie par
$\lforall[x][\lforall[y][(\Atom{Q}{x,y} \lif \Atom{Q}{y, x})]]$ :
une !!{structure} du premier ordre~$\Struct{M}$ telle que
$\Domain{M} = W$ et $\Assign{Q}{M} = R$ satisfait cet
énoncé si et seulement si $R$ est symétrique. Cela suggère
la définition suivante.

\begin{defn}
  Une classe~$\mClass{F}$ de cadres est
  \emph{définissable au premier ordre} s'il existe un
  !!{sentence}~$!A$ dans le langage du premier ordre ayant
  pour seul !!{predicate} binaire~$Q$, tel que
  $\mModel{F} = \tuple{W, R} \in \mClass{F}$ si et seulement si
  $\Sat{M}{!A}$ dans la !!{structure} du premier
  ordre~$\Struct{M}$ où $\Domain{M} = W$ et
  $\Assign{Q}{M} = R$.
\end{defn}

Les propriétés étudiées jusqu'ici et les !!{formula}s modales
qui les définissent font toutefois figure d'exceptions. Une
!!{formula} modale ne définit pas nécessairement une classe
de cadres définissable au premier ordre; inversement, toute
classe de cadres définissable au premier ordre n'est pas
nécessairement définissable par une formule modale.

La formule de Löb fournit un contre-exemple à la première
affirmation :
\begin{equation}
\Box(\Box p \lif p) \lif \Box p. \tag{\Ax{W}}
\end{equation}
\Ax{W} définit la classe des cadres transitifs dont la relation
est inversement bien fondée. Une relation est \emph{bien fondée}
s'il n'existe aucune suite infinie $w_1$, $w_2$, \dots{} telle que
$Rw_2w_1$, $Rw_3w_2$, \dots. Par exemple, la relation $<$
sur~$\Nat$ est bien fondée, contrairement à $<$ sur~$\Int$.
Une relation est \emph{inversement bien fondée} si sa relation
réciproque est bien fondée. Ainsi, il n'existe pour elle
aucune suite infinie $w_1$, $w_2$, \dots{} telle que
$Rw_1w_2$, $Rw_2w_3$, \dots.

Il n'existe cependant aucun énoncé du premier ordre qui définisse
la classe des relations transitives et inversement bien fondées.
Supposons en effet que~$\Sat{M}{!F}$ si et seulement si
$R = \Assign{Q}{M}$ est transitive et inversement bien fondée.
Étendons le langage par de nouvelles constantes
$a_1$, $a_2$, \dots{} et, pour chaque $n \ge 2$,
définissons l'!!{sentence}~$!A_n$ par
\[
(\Atom{Q}{a_1,a_2} \land \dots \land
    \Atom{Q}{a_{n-1},a_{n}})
\]
Considérons alors l'ensemble d'!!{sentence}s
\[
\Gamma = \{!F, !A_2, !A_3, \dots\}.
\]
Tout sous-ensemble fini de $\Gamma$ est satisfaisable.
Choisissons $k \ge 2$ au moins aussi grand que tous les indices
$i$ pour lesquels $!A_i$ appartient au sous-ensemble considéré;
si aucun n'y figure, prenons $k = 2$. Construisons une
!!{structure}~$\Struct{M_k}$ telle que
$\Domain{M_k} = \{1, \dots, k\}$ et
$\Assign{Q}{M_k} = <$, avec
$\Assign{a_i}{M_k} = i$ pour $i \le k$ et
$\Assign{a_i}{M_k} = 1$ pour $i > k$.
La relation $<$ sur $\{1, \dots, k\}$ est transitive et
inversement bien fondée, donc $\Sat{M_k}{!F}$.
Par construction, $\Sat{M_k}{!A_i}$ pour tout $i$ tel que
$2 \le i \le k$.
Le théorème de compacité du premier ordre donne alors un
modèle~$\Struct{M}$ de $\Gamma$. Puisque
$\Sat{M}{!F}$, la relation $\Assign{Q}{M}$ est
inversement bien fondée par hypothèse. Or, puisque
tous les $!A_n$ appartiennent à $\Gamma$, la suite
$\Assign{a_1}{M}$, $\Assign{a_2}{M}$, \dots{}
forme une chaîne infinie pour $\Assign{Q}{M}$,
contrairement à cette propriété.

La propriété d'universalité fournit un contre-exemple à la
seconde affirmation : $Ruv$ pour tous les mondes $u$ et $v$.
La classe des cadres universels est définissable au premier
ordre par $\lforall[x][\lforall[y][\Atom{Q}{x,y}]]$.
Cependant, aucune formule modale n'est valide dans tous les
cadres universels et dans ceux-là seulement. En effet, les
formules valides dans tous les cadres universels sont
exactement celles qui sont valides dans tous les cadres
réflexifs, symétriques et transitifs. Tout cadre de ce
dernier type se décompose en classes d'équivalence; dans
chacune, la relation est universelle, et la vérité d'une
formule modale en un monde ne dépend que de sa classe.
Or certains cadres réflexifs, symétriques et transitifs ne
sont pas universels. Toute !!{formula} valide dans tous
les cadres universels est donc également valide dans
certains cadres non universels.
Dans la dernière preuve, l'invariance par restriction à une
classe d'équivalence se vérifie par induction sur les formules.
La valuation se restreint à cette classe et, à tout monde de
la classe, l'ensemble des mondes accessibles reste exactement
le même. Les clauses des modalités sont donc préservées.
\end{document}
